% Lösungen Einheit 2 von 5 – Streckenlängen, Mittelpunkt und Teilpunkte (zusammenbau v0.1)
\einheitenkopf{Einheit 2 von 5 · Streckenlängen, Mittelpunkt und Teilpunkte}

\erg{13}{a) ein Punkt – die Mitte von $AB$ \quad b) $\overrightarrow{AB} = (2 | 3 | 6)$, $|\overrightarrow{AB}| = \sqrt{4 + 9 + 36} = 7$ LE \quad c) $|\overrightarrow{SH}| = \sqrt{4 + 9 + 36} = 7$ m; mit Zuschlag $0{,}4$ m (nicht $40$ m): $4 \cdot 7{,}4 = 29{,}6$ m \quad d) $|\overrightarrow{AB}| = \sqrt{36 + 64 + 0} = 10$ m; vorher $10 : 1{,}25 = 8$ m (nicht $10$ minus $25\,\%$) \quad e) $\frac12 \cdot (\overrightarrow{OD} + \overrightarrow{OS}) = (2 | 1{,}5 | 3)$, also ist $W$ Mittelpunkt; an der $x_1x_3$-Ebene gespiegelt: $T(2 | -1{,}5 | 3)$ \quad f) $T_1 = P + \frac{1}{3} \cdot \overrightarrow{PQ} = (6 | 1 | 2)$, $T_2 = P + \frac{2}{3} \cdot \overrightarrow{PQ} = (3 | 2 | 1)$ (nicht die Mitte) \quad g) $Q_h = M + k \cdot \overrightarrow{MD}$ mit $9k = h$: $Q_h(3 - \frac{h}{3} | 3 - \frac{h}{3} | h)$; $Q_{3}(2 | 2 | 3)$ \quad h) $D = A - \overrightarrow{AB} = (1 | 4 | 0)$ (nicht $B$ und nicht $A + 2\overrightarrow{AB}$) \quad i) $\overrightarrow{AC} = \overrightarrow{OC} - \overrightarrow{OA} = 2 \cdot \overrightarrow{OA} = (4 | -2 | 4)$, $|\overrightarrow{AC}| = 6$ (nicht $|\overrightarrow{OC}|$) \quad j) $A$ und $C$ sind Mittelpunkte gegenüberliegender Würfelflächen, ihr Abstand ist die Würfelkante (nicht eine Oktaederkante): $|\overrightarrow{AC}| = \sqrt{16 + 16 + 4} = 6$.}
\erg{14}{a) Fehler: Jan quadriert die Koordinaten von $B$ statt der Differenzen. Richtig: $\overrightarrow{AB} = (2 | 1 | 2)$, $|\overrightarrow{AB}| = \sqrt{4 + 1 + 4} = 3$. \quad b) Pythagoras zweimal: In der Grundebene hat die Diagonale die Länge $\sqrt{a^2 + b^2}$; mit der senkrechten Kante $c$ bildet sie ein rechtwinkliges Dreieck, also ist die Länge $\sqrt{a^2 + b^2 + c^2}$.}
